\section{Discussion}
\label{sec:discussion}

\textbf{Aggregate benchmark scores can be misleading.}
The affordance subset appears substantially easier than spatial reasoning in the headline table, but \cref{tab:subcat} shows that this is not uniform affordance competence. Object blockage/reachability questions drive most of the score, while grasp stability remains below 50\% for every audited run. The bottle scene in \cref{fig:data-samples} illustrates the kind of visual judgment involved: the model must distinguish the target object, the robot arm, and possible obstructions rather than answer from the word ``obstacle'' alone. This matters for robotics because a robot that can answer obstacle questions from template cues may still fail to judge whether a grasp is physically stable.

\textbf{Spatial reasoning remains a critical weakness.}
The audited spatial results cluster around chance despite evaluating models with strong general visual-language capabilities~\cite{bai2025qwen25vl,wu2024deepseekvl2}. Directional orientation, depth/distance ordering, and cross-view correspondence questions all remain weak. The visual categories in \cref{fig:data-samples} show why these questions are nontrivial: colored overlays must be interpreted relative to camera geometry, robot pose, and multi-view alignment. The qualitative evidence in \cref{tab:evidence} shows one concrete failure mode: a CoT response infers camera distance from 2D point position, which is an unreliable shortcut for the tested views.

\textbf{CoT is an ablation, not a default improvement.}
CoT is often treated as a low-cost way to improve reasoning. Here, it is at best mixed. Qwen changes many individual decisions without improving aggregate accuracy. DeepSeek gains a small amount on spatial questions but fails dramatically on affordance CoT, especially object blockage/reachability. The first row of \cref{tab:evidence} shows the failure directly: in a bottle reachability question, zero-shot predicts the correct answer, while CoT treats the target bottle itself as an obstacle and changes to the wrong answer. This is not merely random noise; it matches the aggregate fall from 84.8\% to 13.6\% on object blockage/reachability.

\begin{table*}[t]
  \centering
  \small
  \setlength{\tabcolsep}{4pt}
  \begin{tabular}{@{}p{0.18\linewidth}p{0.30\linewidth}p{0.22\linewidth}p{0.22\linewidth}@{}}
    \toprule
    Claim tested & Example question & Observed outputs & Why it matters \\
    \midrule
    DeepSeek CoT can harm affordance reasoning &
    Is there any obstacle blocking the robot from reaching bottle? &
    Expected A. DeepSeek ZS: A, correct. DeepSeek CoT: D, wrong. CoT excerpt: ``presence of an object (bottle) ... may pose a challenge.'' &
    The reasoning trace confuses the target with an obstacle, explaining the object blockage/reachability collapse. \\
    \midrule
    CoT can use weak 2D shortcuts for depth &
    Which colored point is farthest from the camera? &
    Expected B. Qwen ZS: B, correct. Qwen CoT: D, wrong. CoT excerpt: ``distance ... inferred from their positions.'' &
    The explanation relies on image-plane position as a proxy for 3D distance. \\
    \midrule
    CoT can help individual cases without improving aggregate accuracy &
    Is there any obstacle blocking the robot from reaching chip bag? &
    Expected C. Qwen ZS: D, wrong. Qwen CoT: C, correct. CoT notes no obstacle between robot and chip bag. &
    This explains why per-example behavior changes substantially even when Qwen's aggregate CoT score decreases. \\
    \midrule
    Answer explanations can be ungrounded &
    Wrong-answer probe on a spatial question with correct answer A: Red but proposed answer B: Purple. &
    Qwen response: ``The proposed answer B: Purple is supported by the visible evidence.'' &
    The model can rationalize a deliberately wrong answer instead of rejecting it. \\
    \midrule
    Image mismatch detection differs by model &
    Shuffled-image probe with original correct answer A: Red. &
    Qwen: unclear description of the mismatched scene. DeepSeek: ``current image does not visibly support'' the answer. &
    The two models can have similar spatial accuracy while differing sharply in grounding diagnostics. \\
    \bottomrule
  \end{tabular}
  \caption{Representative evidence from the actual evaluation and sanity-probe artifacts. These examples are not additional metrics; they expose the failure mechanisms behind the aggregate findings.}
  \label{tab:evidence}
\end{table*}

\textbf{Grounded explanations need separate tests.}
The sanity probes reveal a distinction between answering and grounding. Qwen correctly rejects blank images, yet it also supports deliberately wrong answers and often fails to reject shuffled images. DeepSeek handles shuffled images better, but its explanations are often not phrased as explicit support for the known correct answer. The last two rows of \cref{tab:evidence} make this failure visible: the same evaluation setup can expose rationalization, mismatch sensitivity, and answer accuracy as separate properties. A benchmark that reports only multiple-choice accuracy would miss these behaviors.

\textbf{Limitations.}
The main raw-result matrix uses 200 examples per subset and file-order sampling. The sanity probe heuristic is coarse and should be followed by manual review before making fine-grained claims about explanation quality. The benchmark also inherits Robo2VLM's multiple-choice format, template distribution, and generated-answer assumptions. These limitations do not weaken the central diagnostic result: targeted slicing and controls expose failures that aggregate robotics VQA scores can hide.

The scope of the evaluation was constrained by available computational resources, limiting the study to relatively small VLMs. Since model scaling has been shown to improve performance across a range of reasoning tasks, it remains an open question whether larger models would exhibit stronger spatial and affordance reasoning abilities. Exploring this relationship constitutes an important avenue for future research.

\textbf{Flawed dataset.}
A key finding from the human evaluation is that performance is not consistently high even for human participants. Several samples were identified as overly ambiguous, under-specified, or otherwise too difficult to be reliably solved based on the provided image and question alone. This indicates that a non-negligible portion of the dataset does not constitute well-posed evaluation cases for spatial and affordance reasoning.

As a consequence, the benchmark in its current form cannot be directly used as-is for reliable model evaluation. Instead, it requires additional manual filtering and curation to isolate representative and well-defined task instances. Such handpicking of high-quality samples is necessary to ensure that the final benchmark accurately reflects the intended reasoning capabilities and provides a meaningful basis for comparing model performance.
